\chapter{Μεθοδολογία}
\label{chap:methodology}

\section{Το σύνολο δεδομένων}

Χρησιμοποιείται το σύνολο δεδομένων \textenglish{Online Retail II} του
αποθετηρίου \textenglish{UCI}, το οποίο καταγράφει τις συναλλαγές βρετανικού
ηλεκτρονικού καταστήματος χονδρικής πώλησης ειδών δώρων για την περίοδο
2009-2011.

Η επιλογή του στηρίχθηκε σε τρία κριτήρια. Είναι δημόσιο, επομένως τα
αποτελέσματα είναι αναπαραγώγιμα και ελέγξιμα. Αποτελεί το σύνολο δεδομένων
της εργασίας των \textcite{chen2012data}, γεγονός που επιτρέπει τη σύγκριση με
δημοσιευμένη έρευνα. Τέλος, περιέχει ακυρωμένα παραστατικά, χωρίς τα οποία το
ερώτημα των επιστροφών δεν θα ήταν απαντήσιμο.

\subsection{Δομή των εγγραφών}

Κάθε εγγραφή αντιστοιχεί σε μία γραμμή προϊόντος εντός ενός παραστατικού και
περιλαμβάνει αναγνωριστικό παραστατικού, κωδικό προϊόντος, περιγραφή,
ποσότητα, ημερομηνία, τιμή μονάδας, αναγνωριστικό πελάτη και χώρα.

Τα δύο διαθέσιμα φύλλα δεν χρησιμοποιούν τα ίδια ονόματα στηλών. Ο προσαρμογέας
που υλοποιήθηκε δέχεται και τις δύο ονοματοδοσίες, ώστε οποιαδήποτε έκδοση του
συνόλου δεδομένων να λειτουργεί χωρίς τροποποίηση.

\subsection{Δηλωμένοι περιορισμοί}

Το σύνολο δεδομένων επιβάλλει τρεις περιορισμούς που καθορίζουν τη σχεδίαση και
δηλώνονται ρητά:

\begin{itemize}
  \item Δεν περιέχει κόστος πωληθέντων ούτε δαπάνες marketing.
  \item Δεν περιέχει δομημένη κατηγοριοποίηση προϊόντων, παρά μόνο ελεύθερο
        κείμενο περιγραφής.
  \item Η γεωγραφική πληροφορία περιορίζεται σε επίπεδο χώρας.
\end{itemize}

\section{Αρχιτεκτονική του συστήματος}

Το σύστημα υλοποιήθηκε ως εφαρμογή \textenglish{FastAPI} με βάση δεδομένων
\textenglish{PostgreSQL} και ουρά εργασιών \textenglish{Celery}. Η δομή του
είναι στρωματωμένη: διεπαφή, υπηρεσίες, αποθετήρια και μοντέλα.

Οι αλγόριθμοι εξόρυξης και μηχανικής μάθησης υλοποιούνται σε ξεχωριστό επίπεδο,
χωρίς καμία εξάρτηση από τη βάση δεδομένων ή το πλαίσιο της εφαρμογής. Δέχονται
απλές δομές δεδομένων και επιστρέφουν απλές δομές δεδομένων.

Ο διαχωρισμός αυτός δεν είναι αισθητικός. Επιτρέπει τον έλεγχο των αλγορίθμων
με δεδομένα κατασκευασμένα ώστε το αναμενόμενο αποτέλεσμα να υπολογίζεται με
το χέρι, κάτι που θα ήταν ανέφικτο αν οι αλγόριθμοι διάβαζαν απευθείας από τη
βάση, και καθιστά τον ίδιο κώδικα άμεσα αξιοποιήσιμο από σημειωματάρια
\textenglish{Jupyter}.

Οι χρονοβόρες αναλύσεις εκτελούνται ασύγχρονα σε εργάτες, καθώς η εξόρυξη
συχνών υποσυνόλων σε πλήρες ιστορικό λιανικής υπερβαίνει κατά πολύ τη διάρκεια
ζωής ενός αιτήματος.

\section{Κανονικοποιημένο σχήμα δεδομένων}

Οι αναλύσεις δεν διαβάζουν ποτέ το σχήμα μιας συγκεκριμένης πηγής. Κάθε πηγή
αντιστοιχίζεται σε κοινό σχήμα έξι οντοτήτων: πελάτης, προϊόν, περιοχή,
παραγγελία, γραμμή παραγγελίας και επιστροφή.

Τρεις σχεδιαστικές αποφάσεις αξίζουν αναφοράς.

\subsection{Η περιοχή αποθηκεύεται στην παραγγελία}

Η περιοχή στην οποία πραγματοποιήθηκε μια πώληση αποτελεί ιδιότητα της πώλησης
και όχι μόνο του πελάτη. Η αποθήκευσή της στην παραγγελία διατηρεί τα
αποτελέσματα ορθά ακόμη και αν η περιοχή ενός πελάτη διορθωθεί εκ των υστέρων,
και απαλλάσσει από μια σύζευξη το συχνότερο ερώτημα της εργασίας.

\subsection{Οι επιστροφές είναι αυτοτελείς εγγραφές}

Οι επιστροφές καταγράφονται ως ξεχωριστές εγγραφές με θετικά μεγέθη, και όχι ως
αρνητικές γραμμές παραγγελίας.

Πρόκειται για την πλέον καθοριστική απόφαση μοντελοποίησης της εργασίας:
καθιστά δομικά αδύνατη τη σιωπηρή αφαίρεση των επιστροφών από τα έσοδα. Έσοδα
και επιστροφές παραμένουν ανεξάρτητα μετρήσιμα σε κάθε ανάλυση.

\subsection{Η τιμή αποθηκεύεται ανά γραμμή}

Το ίδιο προϊόν πωλείται σε διαφορετικές τιμές στο πέρασμα του χρόνου. Η τιμή
που χρεώθηκε αποθηκεύεται στη γραμμή της παραγγελίας, και η κατηγοριοποίηση ανά
εύρος τιμής εφαρμόζεται σε αυτήν, ώστε κάθε πώληση να προσμετράται εκεί όπου
πράγματι πραγματοποιήθηκε.

\section{Προσαρμογείς πηγών δεδομένων}

Η διεπαφή ενός προσαρμογέα συνίσταται σε μία μέθοδο, η οποία παράγει ροή
κανονικοποιημένων εγγραφών. Η προσάρτηση νέου καταστήματος απαιτεί τη συγγραφή
ενός τέτοιου προσαρμογέα και καμία απολύτως αλλαγή στον αναλυτικό κώδικα.

Οι προσαρμογείς είναι σκοπίμως απλοϊκοί: μεταφέρουν το πρόσημο της ποσότητας
όπως το βρίσκουν, χωρίς να ερμηνεύουν τι σημαίνει. Η ερμηνεία ανήκει στην
υπηρεσία εισαγωγής, ώστε να εφαρμόζεται ταυτόσημα σε κάθε πηγή.

\section{Προεπεξεργασία}

Εγγραφές απορρίπτονται, με καταγραφή του λόγου και του πλήθους, όταν δεν
μπορούν να στηρίξουν ανάλυση: μη αναγνώσιμη ημερομηνία ή ποσότητα, απουσία
κωδικού προϊόντος, μηδενική ποσότητα, μη θετική τιμή.

Η απόρριψη των μη θετικών τιμών χρήζει αιτιολόγησης. Στο συγκεκριμένο σύνολο
δεδομένων οι εγγραφές αυτές αντιστοιχούν σε προσαρμογές αποθέματος και
διαγραφές επισφαλών απαιτήσεων, όχι σε εμπορικές συναλλαγές, και η συμπερίληψή
τους θα στρέβλωνε κάθε μέγεθος εσόδων.

Αντιθέτως, οι εγγραφές χωρίς αναγνωριστικό πελάτη \emph{διατηρούνται}. Αποτελούν
σημαντικό ποσοστό του συνόλου και παραμένουν πραγματικές πωλήσεις. Εξαιρούνται
μόνο από τις αναλύσεις που απαιτούν απόδοση σε πρόσωπο, και κάθε φορά ρητά.

\section{Επαλήθευση της υλοποίησης}

Η ορθότητα των υπολογισμών δεν τεκμηριώνεται με επίκληση της υλοποίησης αλλά
με αυτοματοποιημένους ελέγχους.

Όπου το αποτέλεσμα είναι αριθμητικά προσδιορισμένο, ελέγχεται η ακριβής τιμή
του επί δεδομένων υπολογίσιμων με το χέρι. Όπου υφίσταται μαθηματική ταυτότητα,
όπως στη λογιστική ανάπτυξης, η ταυτότητα επαληθεύεται άμεσα, τόσο σε συνθετικά
όσο και σε πραγματικά δεδομένα. Όπου μια ιδιότητα είναι κρίσιμη, όπως η
τυποποίηση πριν τη συσταδοποίηση, κατασκευάζεται έλεγχος που αποτυγχάνει αν η
ιδιότητα αφαιρεθεί.
